\setcounter{subsection}{7}
\subsection{}
\label{X.2.8}
La última afirmación
\ifthenelse{\boolean{orig}}
{\Ref{X.2.6}}
{de~\Ref{X.2.6}}
deja de ser válida en característica~$p>0$, como se ve ya en el caso de las
curvas elípticas. Como se ha señalado anteriormente, no sabemos si el grupo
fundamental de~$X_0$ es topológicamente de presentación finita; parece por
completo improbable.

\setcounter{proposition}{8}
\begin{theoreme}
\label{X.2.9}
Sean~$k$ un cuerpo algebraicamente cerrado y~$X$ un esquema propio y conexo
sobre~$k$. Entonces el grupo fundamental de~$X$ es topológicamente de
generación finita.
\end{theoreme}

Procedemos
\marginpar{274}
por inducción sobre~$n=\dim X$; la afirmación es trivial para~$n\le 0$.
Supongamos, pues, que~$n>0$ y que el teorema está demostrado en dimensiones
~$n'<n$. Por el lema de Chow (EGA~II 5.6.2), existe un esquema proyectivo~$X'$
sobre~$k$ y un morfismo sobreyectivo~$X'\to X$. Puede suponerse claramente
que~$X'$ es reducido y, pasando a su normalización, normal. Gracias a la teoría
del descenso, basta probar que los grupos fundamentales de las componentes
conexas de~$X'$ son topológicamente de generación finita
(\Ref{IX.5.2}). Quedamos reducidos así al caso en que~$X$ es \emph{proyectivo}
y \emph{normal}. Si~$n=1$, basta aplicar~\Ref{X.2.6}. Si~$n\ge 2$, considérese
una inmersión proyectiva~$X\to\PP_k^r$ y una sección de hiperplano~$Y=X\cdot H$
\ifthenelse{\boolean{orig}}
{(dotado}
{(dotada}
de la estructura reducida inducida), tal que
\ifthenelse{\boolean{orig}}
{$Y\ne X$}
{$Y\ne X$,}
\ie $H\not\supset X$. Se tiene entonces~$\dim Y<n$ y, teniendo en cuenta la
hipótesis de inducción, basta probar que~$\pi_1(Y)\to\pi_1(X)$ es
\emph{sobreyectivo}. Más generalmente:

\begin{lemme}
\label{X.2.10}
Sean~$X$ un preesquema propio sobre un cuerpo algebraicamente cerrado~$k$ y
~$g\colon X\to\PP_k^r$ un morfismo. Supongamos que~$X$ es irreducible y normal
y que~$\dim g(X)\ge 2$. Sean~$H$ un hiperplano de~$\PP_k^r$ y
~$Y=X\times_{\PP_k^r}H$. Entonces~$Y$ es conexo y el homomorfismo
~$\pi_1(Y)\to\pi_1(X)$ es sobreyectivo.
\end{lemme}

Estas afirmaciones se deducen, en efecto, de la siguiente:

\begin{corollaire}
\label{X.2.11}
Bajo las hipótesis anteriores, sean~$X'$ un recubrimiento étale conexo de~$X$ y
~$Y'=X'\times_XY=X'\times_{\PP_k^r}H$ el recubrimiento inducido sobre~$Y$.
Entonces~$Y'$ es conexo.
\end{corollaire}

Como~$X$ es normal,~$X'$ es normal y, por ser conexo,~$X'$ es irreducible;
además, su imagen en~$\PP_k^r$ tiene dimensión~$\ge 2$. Un conocido lema de
Zariski (llamado «\emph{teorema de Bertini}») implica, por tanto, que si
~$H_1'$ es el hiperplano genérico de~$\PP_k^r$, definido sobre una extensión
~$K$ de~$k$, entonces
\ifthenelse{\boolean{orig}}
{$X'\times_{\PP^r}H_1$}
{$X'\times_{\PP^r}H_1'$}
es universalmente irreducible y, por tanto, universalmente conexo sobre~$K$.
El teorema de conexidad de Zariski (EGA~III 4) implica entonces que para
\emph{todo} hiperplano~$H$ (definido sobre una extensión cualquiera de~$k$),
~$X'\times_{\PP^r}H$ es geométricamente conexo. Esto termina la demostración
de~\Ref{X.2.11} y, por consiguiente, la de~\Ref{X.2.9}.

\begin{corollaire}[Lang--Serre]
\label{X.2.12}
Bajo
\marginpar{275}
las hipótesis de~\Ref{X.2.9}, para todo grupo finito~$G$, el conjunto de clases,
salvo isomorfismo, de recubrimientos principales de~$X$ con grupo~$G$ es finito.
\end{corollaire}

\begin{remarque}
\label{X.2.13}
Bajo las hipótesis de~\Ref{X.2.10}, cuando~$\dim g(X)\ge 3$ (al menos cuando
~$g$ es una inmersión y~$X$ es regular), probaremos un resultado más preciso,
conocido en geometría algebraica con el nombre de «\emph{teorema de Lefschetz}»:
~$\pi_1(Y)\to\pi_1(X)$ \emph{es un isomorfismo}.\footnote{\Cf el seminario
SGA~2 (1962), que siguió a este%
\ifthenelse{\boolean{orig}}
{.}
{, teorema X~3.10.}}
En los casos clásicos hay enunciados análogos para los grupos de homología y
los grupos de homotopía superiores, que tarde o temprano habrán de quedar
abarcados por la geometría algebraica abstracta. Ni siquiera para la
cohomología de Hodge~$\H^p(X,\Omega^q)$ parece haberse estudiado aún la
cuestión; por lo demás, es poco probable que para esta última los teoremas de
Lefschetz subsistan sin cambio en característica~$p>0$.
\end{remarque}

\iforig
\else
\begin{remarqueMR}
\label{X.2.14}
Sea~$R$ un anillo de valoración discreta completo, de cuerpo residual
algebraicamente cerrado~$\kres$ de característica~$p>0$ y cuerpo de
fracciones~$K$, y sea~$Y$ una curva propia, lisa y conexa de género~$g$
sobre~$R$. Se tiene, pues, un morfismo de especialización sobreyectivo
~$\sp : \pi_1(Y_{\overline{K}})\to\pi_1(Y_{\kres})$.
Ya se ha observado que, si~$K$ tiene característica~0,~$\sp$ no es un
isomorfismo cuando~$g\geq1$. Supongamos que~$K$ tiene característica~$p$, de
modo que~$R$ es isomorfo al anillo de series formales~$\kres[[T]]$.

En igual característica~$p>0$, el grupo fundamental no queda determinado por
el género~$g$, como se ve ya en las curvas elípticas, que pueden ser ordinarias
o supersingulares. Citemos el resultado reciente de A\ptbl Tamagawa (todavía
inédito). Si~$G$ es un grupo profinito, se denota por~$G^{\res}$ el grupo
profinito cociente de~$G$ que es límite proyectivo de los cocientes topológicos
finitos resolubles de~$G$.

\begin{theoremstar}[A\ptbl Tamagawa]
Supongamos que~$g\geq2$, que la fibra especial~$Y_{\kres}$ es definible sobre
un cuerpo finito y que el morfismo
~$\sp^{\res}:\pi_1(Y_{\overline{K}})^{\res}\to
\pi_1(Y_{\kres})^{\res}$ inducido por~$\sp$ al pasar al cociente es biyectivo.
Entonces la curva~$Y$ es constante sobre~$R$.
\end{theoremstar}

Obsérvese que el grupo de Galois de~$\overline{K}/K$ es resoluble. El enunciado
anterior también puede expresarse así: supongamos que todo recubrimiento finito
étale~$X_K\to Y_K$ de la fibra genérica~$Y_K$, galoisiano con grupo de Galois
resoluble, tiene reducción potencialmente buena (es decir, se extiende a un
recubrimiento finito étale de~$Y$ después de sustituir eventualmente~$R$ por
su normalización en una extensión finita de~$K$). Entonces~$Y$ es constante
sobre~$R$.
\end{remarqueMR}
\fi
